% GENERATED FILE. Edit canonical lessons, then run npm run generate:books.
% canonical-source-hash: fnv1a64:86c496d07972a319
% canonical-lessons: FR-C379-diplome, FR-C379-epreuve, FR-C379-mention, FR-C379-doctorat, FR-C379-bourse

\chapter{Le diplôme}
\label{ch:fr-le-diplome}
\hlchaptermodality{379}

\begin{chapteropening}
\textbf{By the end of this chapter:} \emph{I can talk about qualifications: a diploma, an exam paper, a grade with honours, a doctorate and a grant.}

Complete the last lesson of chapter 379: I can talk about qualifications: a diploma, an exam paper, a grade with honours, a doctorate and a grant.
\end{chapteropening}

\section[le diplôme]{le diplôme — the diploma, the degree}

\label{lesson:FR-C379-diplome}



\begin{quote}
Before the new one: say the French for the average, then the French for the resit.
\end{quote}

\subsection*{You'll want to know: le diplôme}
\textbf{le diplôme} — \textquotedblleft{}the diploma, the degree\textquotedblright{}.

\begin{quote}
Il a enfin son diplôme d'infirmier.
\end{quote}

\begin{quote}
\emph{He finally has his nursing diploma.}
\end{quote}

\emph{Passer un diplôme} — to take a qualification.

\begin{cousinweb}
\emph{diplôme} comes from Latin \emph{diplōma} \textquotedblleft{}a letter folded double, a state letter of recommendation\textquotedblright{} (EtymDB; Lewis \& Short), from Greek \emph{díplōma} \textquotedblleft{}a folded paper\textquotedblright{}, built on \emph{diplóos} \textquotedblleft{}double\textquotedblright{}.
\end{cousinweb}

\begin{grammarlens}[title={What you've built}]
One more word for telling, explaining and describing at B1.
\end{grammarlens}

\subsection*{Your turn}
\begin{itemize}
\raggedright
  \item \emph{Say it:} \emph{le diplôme}
  \item \emph{Say it:} \emph{le diplôme}, once more
  \item \emph{Say it:} a sentence of your own with \emph{le diplôme}
  \item \emph{Recall:} say the French for the average, then the French for the resit, then say \emph{le diplôme} again
\end{itemize}

\begin{tcolorbox}[breakable,colback=teal!4,colframe=teal!35!black,title={Before you move on}]
What does le diplôme mean? (\textquotedblleft{}The diploma, the degree\textquotedblright{}.) Say it once more.
\end{tcolorbox}
\section[l'épreuve]{l'épreuve — the test, the exam paper; the ordeal}

\label{lesson:FR-C379-epreuve}



\begin{quote}
Before the new one: say the French for the resit, then the French for the diploma.
\end{quote}

\subsection*{You'll want to know: l'épreuve}
\textbf{l'épreuve} — \textquotedblleft{}the test, the exam paper; the ordeal\textquotedblright{}.

\begin{quote}
La première épreuve de l'examen dure trois heures.
\end{quote}

\begin{quote}
\emph{The first paper of the exam lasts three hours.}
\end{quote}

\emph{Une épreuve écrite}, \emph{une épreuve orale}. In sport, an event. Feminine.

\begin{cousinweb}
\emph{épreuve} is formed on \emph{éprouver} \textquotedblleft{}to test, to feel\textquotedblright{} (EtymDB), built from \emph{é-} and \emph{prouver}, which is inherited from Latin \emph{probāre} \textquotedblleft{}to test, to examine\textquotedblright{} (Lewis \& Short).
\end{cousinweb}

\begin{grammarlens}[title={What you've built}]
One more word for telling, explaining and describing at B1.
\end{grammarlens}

\subsection*{Your turn}
\begin{itemize}
\raggedright
  \item \emph{Say it:} \emph{l'épreuve}
  \item \emph{Say it:} \emph{l'épreuve}, once more
  \item \emph{Say it:} a sentence of your own with \emph{l'épreuve}
  \item \emph{Recall:} say the French for the resit, then the French for the diploma, then say \emph{l'épreuve} again
\end{itemize}

\begin{tcolorbox}[breakable,colback=teal!4,colframe=teal!35!black,title={Before you move on}]
What does l'épreuve mean? (\textquotedblleft{}The test, the exam paper; the ordeal\textquotedblright{}.) Say it once more.
\end{tcolorbox}
\section[la mention]{la mention — the distinction, the honours (on an exam); the mention}

\label{lesson:FR-C379-mention}



\begin{quote}
Before the new one: say the French for the diploma, then the French for the test.
\end{quote}

\subsection*{You'll want to know: la mention}
\textbf{la mention} — \textquotedblleft{}the distinction, the honours (on an exam); the mention\textquotedblright{}.

\begin{quote}
Elle a eu son diplôme avec mention.
\end{quote}

\begin{quote}
\emph{She got her diploma with honours.}
\end{quote}

\emph{Mention bien}, \emph{mention très bien} — with honours, with high honours.

\begin{cousinweb}
\emph{mention} is a learned borrowing from Latin \emph{mentiō} \textquotedblleft{}a calling to mind, a mention\textquotedblright{} (EtymDB; Lewis \& Short), related to \emph{mēns} \textquotedblleft{}mind\textquotedblright{}.
\end{cousinweb}

\begin{grammarlens}[title={What you've built}]
One more word for telling, explaining and describing at B1.
\end{grammarlens}

\subsection*{Your turn}
\begin{itemize}
\raggedright
  \item \emph{Say it:} \emph{la mention}
  \item \emph{Say it:} \emph{la mention}, once more
  \item \emph{Say it:} a sentence of your own with \emph{la mention}
  \item \emph{Recall:} say the French for the diploma, then the French for the test, then say \emph{la mention} again
\end{itemize}

\begin{tcolorbox}[breakable,colback=teal!4,colframe=teal!35!black,title={Before you move on}]
What does la mention mean? (\textquotedblleft{}The distinction, the honours (on an exam); the mention\textquotedblright{}.) Say it once more.
\end{tcolorbox}
\section[le doctorat]{le doctorat — the doctorate, the PhD}

\label{lesson:FR-C379-doctorat}



\begin{quote}
Before the new one: say the French for the test, then the French for the distinction.
\end{quote}

\subsection*{You'll want to know: le doctorat}
\textbf{le doctorat} — \textquotedblleft{}the doctorate, the PhD\textquotedblright{}.

\begin{quote}
Mon cousin fait un doctorat en chimie.
\end{quote}

\begin{quote}
\emph{My cousin is doing a PhD in chemistry.}
\end{quote}

The holder is \emph{un docteur}.

\begin{cousinweb}
\emph{doctorat} comes from Medieval Latin \emph{doctōrātus}, built on \emph{doctor} \textquotedblleft{}teacher\textquotedblright{}, from \emph{docēre} \textquotedblleft{}to teach\textquotedblright{} (Lewis \& Short).
\end{cousinweb}

\begin{grammarlens}[title={What you've built}]
One more word for telling, explaining and describing at B1.
\end{grammarlens}

\subsection*{Your turn}
\begin{itemize}
\raggedright
  \item \emph{Say it:} \emph{le doctorat}
  \item \emph{Say it:} \emph{le doctorat}, once more
  \item \emph{Say it:} a sentence of your own with \emph{le doctorat}
  \item \emph{Recall:} say the French for the test, then the French for the distinction, then say \emph{le doctorat} again
\end{itemize}

\begin{tcolorbox}[breakable,colback=teal!4,colframe=teal!35!black,title={Before you move on}]
What does le doctorat mean? (\textquotedblleft{}The doctorate, the PhD\textquotedblright{}.) Say it once more.
\end{tcolorbox}
\section[la bourse]{la bourse — the grant, the scholarship; the stock exchange}

\label{lesson:FR-C379-bourse}



\begin{quote}
Before the new one: say the French for the distinction, then the French for the doctorate.
\end{quote}

\subsection*{You'll want to know: la bourse}
\textbf{la bourse} — \textquotedblleft{}the grant, the scholarship; the stock exchange\textquotedblright{}.

\begin{quote}
Elle a reçu une bourse pour étudier à l'étranger.
\end{quote}

\begin{quote}
\emph{She got a grant to study abroad.}
\end{quote}

\emph{La Bourse} is the stock exchange. Once a purse.

\begin{cousinweb}
\emph{bourse} is inherited through Old French \emph{borse} \textquotedblleft{}purse\textquotedblright{} (EtymDB) from Late Latin \emph{bursa} \textquotedblleft{}a leather bag\textquotedblright{}, from Greek \emph{búrsa} \textquotedblleft{}a hide\textquotedblright{}.
\end{cousinweb}

\begin{grammarlens}[title={What you've built}]
One more word for telling, explaining and describing at B1.
\end{grammarlens}

\subsection*{Your turn}
\begin{itemize}
\raggedright
  \item \emph{Say it:} \emph{la bourse}
  \item \emph{Say it:} \emph{la bourse}, once more
  \item \emph{Say it:} a sentence of your own with \emph{la bourse}
  \item \emph{Recall:} say the French for the distinction, then the French for the doctorate, then say \emph{la bourse} again
\end{itemize}

\begin{tcolorbox}[breakable,colback=teal!4,colframe=teal!35!black,title={Before you move on}]
What does la bourse mean? (\textquotedblleft{}The grant, the scholarship; the stock exchange\textquotedblright{}.) Say it once more.
\end{tcolorbox}
